\section{Trade-offs, Limitations and Scope}\label{sec:tradeoffs}

\subsection{Trade-offs}
Table~\ref{tab:tradeoffs} collects the design trade-offs of the receiver, each quantified by an
experiment of Section~\ref{sec:results}.

\begin{table}[!htbp]\centering
\caption{Design trade-offs and their measured consequences.}\label{tab:tradeoffs}
\resizebox{\textwidth}{!}{%
\begin{tabular}{p{3.6cm}p{6.2cm}p{6.6cm}}\toprule
Choice & Gain & Cost (measured) \\\midrule
Pilot symbols per refresh, $N_p$ & Accuracy: TDL 100\,ns gap 0.7 $\to$ 0.1--0.24 for $N_p=2\to8$ & Overhead $N_p/N$; for fast users one pilot refreshed often is better (net rate, Table~\ref{tab:overhead}) \\
Refresh interval $N$ & Lower overhead & Estimate ages: 120\,km/h needs $N\approx8$ on TDL (12\%); without phase tracking $N\approx4$ (25\%) \\
Delay window $\tau_{\max}$, basis size $L$ & Covers longer echoes (TDL 100\,ns needs 1\,$\mu$s) & More unknowns $\Rightarrow$ more pilots; window beyond the CP is useless (TDL-D, 300\,ns) \\
Rank $r$ & Captures more paths & Little effect between 8 and 32 on TDL; $r=6$ instead of 8 is free and 1.7$\times$ cheaper \\
Distinct patterns $G$ (security, scrambling) & Stronger scrambling for others & Unknowns $+Gr^2$: gap 0.10 $\to$ 1.03 at 30\,dB for $G=1\to16$ with 2 pilots; 4 pilots restore 0.38 \\
Duty cycle $D$ & Away from 0.5: more energy on $h=0$, better one-tap & Our receiver is indifferent (gap 0.05--0.23 for $D=1/8..7/8$); the one-tap ceiling $(2D-1)^2/(1-(2D-1)^2)$ changes \\
Decision-directed passes & 28\,dB/64-QAM: SER ratio 1.54 $\to$ 1.14 with one pass & A second pass with 64-QAM raised SER to 0.24 (error propagation); we use at most one \\
Evidence timing search & Removes the timing requirement; no rate loss vs known timing & 55 first-stage evaluations per search, about $1.2\cdot10^{12}$ operations at $M=3276$ (computed, not yet measured); needed only when the offset changes \\
Two-timescale learning & fewer per-frame unknowns (gain shown only at 300\,ns, a CP-violating case) & Needs $\sim$100 frames sharing the BS--RIS hop; stale if the BS or RIS moves \\
Warm start & 12$\times$ fewer operations (8.6$\times$ faster at full band) & +0.03\,bit/s/Hz gap \\
Freezing $\bm V$ in updates & $\sim$100$\times$ fewer operations & +0.23\,bit/s/Hz at 30\,dB unless the input side is learned with $r=Q$ \\
CG tolerance / iterations & Linear cost in $M$; fits where dense does not & Loose tolerance ($10^{-2}$) left SER 1.18$\times$ dense at $M=2048$; iterations grow with bandwidth \\
Switching rise time & Lower emission (ACLR 36.8 $\to$ 52.7\,dB for 0 $\to$ 0.05 slot) & Shrinks our advantage over truncated estimators (CE-BEM 4.7 $\to$ 8.2 at 0.3 slot); calibrated kernels needed above $\sim$40\,ns \\
Comb pilots & Same accuracy as full pilot symbols, NR-compatible & Neighbouring resource elements must be empty (data next to pilots: 9.79 $\to$ 5.10 at 30\,dB) \\\bottomrule
\end{tabular}}
\end{table}

\subsection{Limitations and scope}
\begin{itemize}
\item \emph{Simulation.} All results are simulated with one published element model; mutual
coupling between elements is ignored, as in the anchor model. Over-the-air validation is future work.
\item \emph{Pattern knowledge.} The receiver must know every panel's switching pattern (A1); with a
wrong pattern for one of two panels the rate collapses from 5.5 to 1.2\,bit/s/Hz.
\item \emph{Delay spread beyond the prefix.} Channels whose taps exceed the cyclic prefix
(TDL-D at 300\,ns) are out of scope of any prefix-limited window; 300\,ns TDL-A/C need the static
side learned over $\sim$100 frames.
\item \emph{High speed.} At 120\,km/h a pilot symbol every $\sim$8 symbols is needed on TDL
channels (12\% overhead at 20\,dB). Two tracking receivers did not improve on this, and at 30\,dB the
in-symbol Doppler mixing, not modelled by our estimator, limits every receiver. A Doppler-aware
extension of the model is future work.
\item \emph{Mobility metric.} All mobility numbers assume common-phase tracking per symbol, as in
5G NR; without it pilots are needed about twice as often.
\item \emph{Hardware.} Switching jitter must stay below about 4\,ns at 30\,dB, residual CFO below
about 1\% of the spacing, and oscillator linewidth near 100\,Hz or tracked; transients are handled by
fitting the rise time.
\item \emph{Low SNR.} At 10\,dB the gain over the one-tap receiver for a user the pattern is designed
for is modest (+0.6\,bit/s/Hz, about 2\,dB): noise, not coupling, dominates there.
\item \emph{Motivation and secrecy.} For a single user an optimized static RIS gives more rate than a TM-RIS, also with
our receiver (Section~\ref{sec:res-motiv}); our receiver reduces the cost of time modulation, it does not make it profitable.
An eavesdropper that learns the pattern can run the same receiver (secrecy rate 0.5--2.9 instead of 1.0--3.6\,bit/s/Hz at 20\,dB).
\item \emph{Rank.} With $Q=16$ and $r=8$ (or $r=Q$ in Algorithm~\ref{alg:learn}) the low-rank property
is weak; it matters for large surfaces, where it is tested only up to $Q=1024$ with narrowband array phases.
\item \emph{Equalizer speed.} The measured CG advantage counts the dense factorization against 12
symbols; with more symbols per estimate the dense factorization is amortized, and the break-even
number of symbols has not been measured.
\item \emph{Statistics.} Some points use 4--8 channels (300\,ns, TDL overhead) and the full-band
28\,dB result 10 channels; their confidence intervals are wider than those of the 20--30-channel results.
\item \emph{Learning-based estimators} were not compared (Section~\ref{sec:method}).
\end{itemize}
